\documentclass[12pt,a4paper]{article}
% ============================================================
%  Structural Persistence Series — Shared Preamble
% ============================================================
\usepackage{fontspec}
\setmainfont{Hiragino Mincho ProN}
\setsansfont{Hiragino Sans}
\setmonofont{Menlo}

\usepackage{iftex}
\ifXeTeX
  \XeTeXlinebreaklocale "ja"
  \XeTeXlinebreakskip=0pt plus 1pt minus 0.1pt
\fi

\usepackage[margin=22mm]{geometry}
\usepackage{amsmath,amssymb,amsthm}
\usepackage{xcolor}
\usepackage{graphicx}
\usepackage{hyperref}
\usepackage{booktabs}
\usepackage{fancyhdr}
% enumitem not available in this TeX installation

% --- Colours ------------------------------------------------
\definecolor{PaperAccent}{HTML}{1F4E79}
\definecolor{PaperGray}{HTML}{51606F}

\hypersetup{
  colorlinks=true,
  linkcolor=PaperAccent,
  urlcolor=PaperAccent,
  citecolor=PaperAccent,
  pdflang=ja
}
\urlstyle{same}

% --- Typography ---------------------------------------------
\setlength{\parindent}{1.5em}
\setlength{\parskip}{0pt}
\setlength{\emergencystretch}{3em}
\linespread{1.08}
\setlength{\abovedisplayskip}{8pt plus 2pt minus 4pt}
\setlength{\belowdisplayskip}{8pt plus 2pt minus 4pt}
\setlength{\abovedisplayshortskip}{6pt plus 2pt minus 2pt}
\setlength{\belowdisplayshortskip}{6pt plus 2pt minus 2pt}

% --- Section label names ------------------------------------
\renewcommand{\abstractname}{要旨}

% --- Theorem-like environments ------------------------------
\theoremstyle{plain}
\newtheorem{proposition}{命題}[section]
\newtheorem{lemma}[proposition]{補題}
\newtheorem{corollary}[proposition]{系}

\theoremstyle{definition}
\newtheorem{definition}{定義}[section]
\newtheorem{assumption}{条件}[section]

\theoremstyle{remark}
\newtheorem*{remark}{注}

% --- Running header -----------------------------------------
\pagestyle{fancy}
\fancyhf{}
\renewcommand{\headrulewidth}{0.4pt}
\fancyhead[L]{\small\textcolor{PaperGray}{\nouppercase{\leftmark}}}
\fancyhead[R]{\small\textcolor{PaperGray}{\thepage}}
\fancypagestyle{plain}{%
  \fancyhf{}
  \renewcommand{\headrulewidth}{0pt}
  \fancyfoot[C]{\small\textcolor{PaperGray}{\thepage}}
}

% --- Title block macro --------------------------------------
\newcommand{\PaperTitleBlock}[3]{%
  % #1 = title, #2 = subtitle, #3 = date
  \thispagestyle{plain}%
  \begin{center}
  {\LARGE\bfseries #1\par}
  \vspace{0.4em}
  {\normalsize #2\par}
  \vspace{1.0em}
  {\large Akihito Sunagawa\par}
  \vspace{0.3em}
  {\normalsize #3\par}
  \end{center}
  \vspace{1.6em}
}

% Legacy 2-arg version (backward compat)
\newcommand{\PaperTitleBlockSimple}[2]{%
  \PaperTitleBlock{#1}{}{#2}
}


\begin{document}

\PaperTitleBlock
  {計算コストの構造的予測}
  {— ランダム3-SATにおける存在と発見の分離 —}
  {2026年4月12日}

\begin{abstract}
解が存在するかどうかと、アルゴリズムがそれを発見できるかどうかは、異なる機構に支配される別の問いである。ランダム 3-SAT では、第一モーメント法により

\[
E[\#\text{SAT}] = N_{\text{eff}} \cdot e^{-L}
\]

が成り立つ。ここで \(L = \alpha\) n\textbar ln(7/8)\textbar{} は節密度の再パラメータ化であり、第一モーメント法の指数そのものである。これは解の存在を特徴づけるが、発見については何も言わない。

本稿では計算予算 μ（ソルバーの操作回数）を導入し、条件つき成功確率 P(found\textbar SAT) が μ/μ\_c(L) の関数として乗法的に分離できるかを問う。3種のソルバー族（CDCL、WalkSAT、ランダム探索）を用い、\(n = 100-300\)変数、\(\alpha = 3.0-4.0\)の範囲で実験した結果、三つの知見を得る。

第一に、中央値実行時間は \(\mu_{c}\) ∝ e\^{}\{cL\} としてスケールし、感度指数 c はソルバーに依存する。CDCL では \(c ≈ 0.24\)、WalkSAT では \(c ≈ 0.21\)、ランダム探索では c = 1.0（厳密）である。

第二に、比 μ/μ\_c は有効な正規化変数であり、両ソルバーとも \(\mu/\mu_{c} = 1\)の付近で単調な遷移を示す。ただし遷移の形状は 14-15\% 異なる。

第三に、c の順序 \(c_{CDCL} > c_{WalkSAT}\)は、テストした全問題サイズ（\(n = 100-300\)）で頑健に成立する。

したがって、同一の構造的パラメータ L が存在（\(c = 1\)、数学的恒等式）と発見（\(c < 1\)、アルゴリズム依存）の双方を支配するが、感度は質的に異なる。とくに c は効率ではなく感度を測る量であり、WalkSAT は c が最小であるにもかかわらず絶対コストは最大である。これは、低感度が構造の活用ではなく構造の無視から生じうることを意味する。
\end{abstract}

\section{はじめに}\label{ux306fux3058ux3081ux306b}

ランダム k-SAT の充足可能性閾値は、組合せ論と統計物理学で最も研究されている現象の一つである。空洞法とその厳密化は、解が存在しなくなる臨界節密度 \(\alpha_{c}\) を漸近的精度で決定する。これらは存在の問いに答える。すなわち、与えられた節変数比 α のもとで、充足割当ては高確率で存在するか。

補完的な問いは、はるかに少ない形式的注目しか受けていない。解が存在するとして、それを発見するにはどれだけの計算が必要か。

本稿はこの間隙を、恒等式としては自明だが定量的プログラムとしては非自明な分解によって形式化する。

\[
P(found) = P(SAT) \times P(found|SAT)
\]

左の因子は第一モーメント法でよく特徴づけられている。右の因子------発見確率------が本稿の測定対象である。

中心的知見は、第一モーメント法の指数と同じ構造的パラメータ L が両方の因子を支配するが、その効き方は異なるということである。

\begin{itemize}
\tightlist
\item
  存在: \(E[\#\text{SAT}] = N_{\text{eff}} \cdot e^{-L}\)、\(c = 1\)（数学的恒等式）
\item
  発見: 中央値計算コストは \(\mu_c \propto e^{cL}\)、\(c < 1\)（ソルバー依存）
\end{itemize}

感度指数 c は、ソルバーのコストが構造的複雑性 L の増加にどれだけ強く応答するかを測る。これは効率そのものの尺度ではない。CDCL（\(c ≈ 0.24\)）は WalkSAT（\(c ≈ 0.21\)）より L への感度が高いが、絶対コストでは 10-50 倍速い。したがって、アルゴリズムの洗練は、単にコストを下げるだけでなく、L に対する応答の仕方そのものを変える。

\section{枠組み}\label{ux67a0ux7d44ux307f}

\subsection{存在: 第一モーメント法}\label{ux5b58ux5728-ux7b2cux4e00ux30e2ux30fcux30e1ux30f3ux30c8ux6cd5}

n 変数、\(m = \alpha n\)節のランダム 3-SAT インスタンスを考える。一様ランダム割当てが単一の節を充足する確率は 1 - 2\^{}\{-k\} = 7/8（\(k = 3\)）。充足割当ての期待数は

\[
E[\#\text{SAT}] = 2^n \cdot (7/8)^m = \exp(n \ln 2 - L)
\]

ここで

\[
L ≡ \alpha n|\ln(7/8)| = m · |\ln(7/8)|
\]

は節数 m（あるいは、n を固定したときの節密度 α）の再パラメータ化であり、1節あたり \textbar ln(7/8)\textbar{} ≈ 0.134 ナットの情報損失に対応する。最小形式の累積損失 L と同一の役割を果たす量であり、各節が充足割当ての集合を (7/8) の比率で縮小する過程の対数的累積である。

L は新しい量ではない。第一モーメント法の指数そのものであり、インスタンスの仕様から計算可能である。新しいのは、L が計算コストをも支配するという経験的観測である。

\subsection{発見: 計算予算の導入}\label{ux767aux898b-ux8a08ux7b97ux4e88ux7b97ux306eux5c0eux5165}

μ をソルバーに与えられた計算予算とする（CDCL では衝突数、WalkSAT ではフリップ数、ランダム探索では試行数）。発見確率を

\begin{definition}[発見確率 $S_{\text{solve}}(L, \mu)$]
$$
S_{\text{solve}}(L, \mu) \equiv P(\text{found}|\text{SAT})
$$

これは、インスタンスが充足可能であることを条件として、ソルバーが予算 $\mu$ 以内に充足割当てを見つける確率である。
\end{definition}

本稿の仮説は、\(S_{solve}\) が中央値実行時間 μ\_c(L) を通じた乗法的分離を許すことである。

\[
S_{solve}(L, \mu) = f(\mu/\mu_{c}(L))
\]

ここで f は単調増加関数であり、中央値実行時間は指数的にスケールする。

\[
\mu_{c}(L) = A · e^{cL}
\]

c が感度指数------構造的複雑性 L の単位あたりの計算コスト成長率------であり、A は実装依存の切片である。

\begin{definition}[指数形の根拠]
制約がコストに独立に寄与するならば、すなわち $\mu_c(L_1 + L_2) = \mu_c(L_1) \cdot [\mu_c(L_2)/A]$ が成り立つならば、コーシーの関数方程式の一意解として $e^{cL}$ が得られる。したがって、この指数形は独立蓄積と両立する唯一の連続的スケーリング則である。
\end{definition}

\begin{definition}[$c = 1$ の基準線]
ランダム探索（一様に割当てを抽出して検査）の場合、解を見つけるまでの期待試行数は $2^n / E[\#\text{SAT}] = e^L$。よって $\mu_c^{\text{random}} \propto e^L$、すなわち $c = 1$ が厳密に成り立つ。$c < 1$ のソルバーは、構造なし基準線より構造的複雑性あたりのコスト成長が小さい。
\end{definition}

\section{実験設定}\label{ux5b9fux9a13ux8a2dux5b9a}

\subsection{インスタンス生成}\label{ux30a4ux30f3ux30b9ux30bfux30f3ux30b9ux751fux6210}

n ∈ \{100, 200, 300\} 変数、節密度 α ∈ \{3.0, 3.5, 3.86, 4.0\} のランダム 3-SAT インスタンスを生成する。\(\alpha > 4.2\)は除外する。充足可能性閾値 \(\alpha_{c} ≈ 4.27\)の近傍では充足可能なインスタンスの割合が急減し、生存者バイアスが生じるためである。各 \((n, \alpha)\) の組に対して 100-200 インスタンスを生成し、充足可能と確認されたもののみを保持する。

\subsection{ソルバー}\label{ux30bdux30ebux30d0ux30fc}

3種のソルバー族が、アルゴリズムの洗練度の幅を覆う。

\begin{definition}[CDCL（PySAT 経由の MiniSat）]
衝突駆動節学習。予算単位は衝突数。
\end{definition}

\begin{definition}[WalkSAT（再起動つき）]
確率的局所探索。予算単位はフリップ数。
\end{definition}

\begin{definition}[ランダム探索]
一様ランダム割当ての検査。$c = 1$ の理論的基準線。
\end{definition}

\section{結果}\label{ux7d50ux679c}

\subsection{感度指数 c}\label{ux611fux5ea6ux6307ux6570-c}

\(\ln(\mu_{c})\) を L に対してプロットし、線形回帰を行うと以下を得る。

{\def\LTcaptype{table} % do not increment counter
\begin{longtable}[]{@{}llll@{}}
\toprule\noalign{}
ソルバー & c & 95\% CI & R² \\
\midrule\noalign{}
\endhead
\bottomrule\noalign{}
\endlastfoot
CDCL（衝突数） & 0.25 & {[}0.24, 0.29{]} & 0.99 \\
WalkSAT（フリップ数） & 0.18 & {[}0.17, 0.22{]} & 0.96 \\
ランダム（試行数） & 1.0 & --- & --- \\
\end{longtable}
}

\(\alpha \ge 3.5\) に限定すると（\(\alpha = 3.0\) でのソルバー起動の床効果を除外）、CDCL で \(c \approx 0.24\)、WalkSAT で \(c \approx 0.21\) となる。

c の順序 \(c_{CDCL} > c_{WalkSAT}\)は当初意外に見える。最も効率的なソルバー（CDCL）が、L への感度がより高い。これは二つの質的に異なる機構を反映している。

\begin{itemize}
\tightlist
\item
  CDCL（\(c \approx 0.24\)）: 節学習が制約構造を能動的に利用する。効率と感度の双方を高める。
\item
  WalkSAT（\(c \approx 0.21\)）: 局所探索は大域的制約構造をほぼ無視する。感度は低いが、絶対コストは 10-50 倍大きい。
\end{itemize}

したがって c と絶対効率（1/A）は独立な2軸であり、完全な特徴づけには両方が要る。

\subsection{μ/μ\_c による正規化}\label{ux3bcux3bc_c-ux306bux3088ux308bux6b63ux898fux5316}

\(S_{solve}\) を \(\ln(\mu/\mu_{c})\) に対してプロットすると、両ソルバーとも \(\mu/\mu_{c} = 1\)の付近で遷移する。シグモイドモデルのフィット結果は以下の通り。

{\def\LTcaptype{table} % do not increment counter
\begin{longtable}[]{@{}llll@{}}
\toprule\noalign{}
フィット & 傾き k & 遷移点 x\_0 & R² \\
\midrule\noalign{}
\endhead
\bottomrule\noalign{}
\endlastfoot
CDCL & 1.36 & -0.04 & 0.985 \\
WalkSAT & 1.57 & 0.00 & 0.972 \\
\end{longtable}
}

両ソルバーとも \(x_0 ≈ 0\)であり、中央値実行時間が計算予算の自然な尺度であることを裏づける。傾き k の差（14\%）は、ワイブル CDF でフィットしても同程度（15\%）であり、遷移形状の差はフィット関数の選択に依存しない。

\section{考察}\label{ux8003ux5bdf}

\subsection{α-n 非対称性}\label{ux3b1-n-ux975eux5bfeux79f0ux6027}

\(L = \alpha\) n\textbar ln(7/8)\textbar{} は α と n に対称的に入るが、計算コストはこの対称性を尊重しない。先行研究により、WalkSAT は n に対して多項式的にスケールする（\(\mu_{c}\) ∼ n\^{}b、\(b ≈ 3\)）が、CDCL は指数的にスケールする（\(\mu_{c}\) ∼ 1.03\^{}n）。

\(\mu_{c} = A\) · e\^{}\{cL\} が両軸で一様に成立するなら、固定 α で \(\mu_{c}\) ∝ e\^{}\{cα\textbar ln(7/8)\textbar·n\} を予測する------任意の \(c > 0\)について n の指数。WalkSAT の多項式 n-スケーリングとは非両立である。

本稿の n-スケーリング実験でもこの非両立を確認した。生存者バイアスのない節密度（\(\alpha = 3.0\), 3.5）で、WalkSAT の中央値フリップ数は n に対して明確に減速する。

{\def\LTcaptype{table} % do not increment counter
\begin{longtable}[]{@{}lll@{}}
\toprule\noalign{}
α & γ（μ\_c ∼ n\^{}γ） & 減速比 \\
\midrule\noalign{}
\endhead
\bottomrule\noalign{}
\endlastfoot
3.0 & 1.07 & 0.48 \\
3.5 & 1.69 & 0.75 \\
\end{longtable}
}

減速比（n: 200→300 の成長率 / 100→200 の成長率）は指数なら 1.0、多項式なら \textless{} 1。両方とも明確に亜指数的である。

CDCL については、切片 \(A(n)\) が n に対して線形に減少し、\(\alpha_{c} ≈ 4.27\)での実効成長率は \(\mu_{c}\) ∝ 1.034\^{}n と予測される。これは先行研究の 1.03\^{}n と整合する。

結論として、L は CDCL のような指数的 n-スケーリングを持つソルバーに対しては有効な統一座標だが、WalkSAT のような多項式的 n-スケーリングを持つソルバーに対しては、c は固定 n での α-感度として解釈すべきである。

\subsection{構造持続の枠組みにおける位置づけ}\label{ux69cbux9020ux6301ux7d9aux306eux67a0ux7d44ux307fux306bux304aux3051ux308bux4f4dux7f6eux3065ux3051}

SAT は、構造を維持できる状態集合（充足割当ての集合）が明示的に定まり、各制約の縮小率も直接計算できるドメインである。累積損失 \(L = \alpha\) n\textbar ln(7/8)\textbar{} はインスタンスの仕様から一意に決定され、代理指標を経由しない。この意味で、SAT は構造持続の形式が最も硬い形で検証されうる地盤を提供する。

この位置づけは、ドメイン汎用フレーム（別稿）における主張強度の段階化と対応する。SAT は Route A（定理ルート）に属する。有効選択肢空間が定義可能であり、各制約の削減率が定義可能であり、削減が加法的情報量に変換可能であり、独立性または弱依存が置ける。存在側の指数表現は第一モーメント法の数学的帰結として成り立ち、発見側の感度指数 c は経験的に測定される。

これに対し、大規模言語モデルにおける論理一貫性維持（別稿）や、前提更新を伴う継続学習（別稿）のような高次元ドメインでは、有効選択肢空間は理論構成物として置かれ、L は代理指標または境界としてしかアクセスできない。これらは Route B（境界ルート）または Route C（実証ルート）に属し、主張強度は弱まる。

本稿群全体では、この強度差を隠すのではなく、明示的に管理する。SAT が与えるのは、L という座標が代理指標なしに機能する最も清潔な事例であり、他ドメインでの主張が Route B/C に弱まることの基準線である。したがって、SAT は構造持続の枠組みの単なる適用例ではなく、その中で最も硬い検証地盤として位置づけられる。

構造持続理論における第一段階の論拠は、構造的パラメータ L が、代理的で単純な基準モデルに対して追加的な予測力を持つことを示す点にある。SAT の文脈で基準線となるのは、ランダム探索の \(c = 1\)である。これは、L ナットあたり一単位のコストを素朴に払う構造非依存の基準線であり、制約の内部構造を一切利用しない。本稿の §4 で示した感度指数の結果は、実用ソルバー（CDCL で \(c ≈ 0.24\)、WalkSAT で \(c ≈ 0.21\)）がこの基準線を大きく下回ることを示している。すなわち、L が定義する制約の内部構造（節学習による将来の縮小の先取り、局所探索の探索地形利用など）を活用すれば、構造非依存の予測からは到達できない効率に達しうる。この結果は、L が単なる制約の個数の再パラメータ化ではなく、構造持続の観点において実質的な予測量であることの直接的な経験的証拠を与える。SAT では L が代理指標を経由せずに定義・測定できるため、この第一段階の論拠は本稿群のどのドメインよりも硬い形で成立する。

ただし、本稿の直接の貢献は、存在と発見の分離、感度指数 c の導入、および α-n 非対称性の明示にあり、構造持続理論の全体を SAT で証明したと主張するものではない。

\section{限界}\label{ux9650ux754c}

本稿の結果は \(n = 100-300\)の範囲に限定されている。\(\alpha \le 3.0\)では CDCL の中央値衝突数が 3-8 と小さく、ソルバー起動の床効果が線形フィットを歪める。\(\alpha > 4.0\)では生存者バイアスが生じる。充足可能性閾値近傍（\(\alpha_{c} ≈ 4.27\)）での c の振る舞いは未解明である。

n-スケーリングの交差予測テストでは、\(n \le 200\) で訓練し \(n = 300\) を予測して、対数空間で \(\text{RMSE} = 0.73\)（約2倍の精度）を達成した。\(\mu_c\) が3桁にわたる範囲での2倍精度は非自明な予測力を示すが、\(n \ge 500\) での検証は未実施である。

SAT 以外のドメインへの拡張には、成果を予測する量とは独立に L を操作的に定義できることが必要であり、その条件は SAT では第一モーメント法により自然に満たされるが、他のドメインでは自明ではない。

\section{結論}\label{ux7d50ux8ad6}

本稿では、ランダム 3-SAT における解の存在と発見を分離した。

\[
P(found) = P(SAT) \times P(found|SAT)
\]

前者は第一モーメント法により \(c = 1\)で支配される。後者------本稿で定量化する対象である------は、中央値計算コストが \(\mu_{c}\) ∝ e\^{}\{cL\}（\(c < 1\)）としてスケールし、μ/μ\_c が発見遷移の有効な正規化変数として機能することを示す。

感度指数 c はソルバーに依存し（CDCL: ≈0.24、WalkSAT: ≈0.21、ランダム: 1.0）、c の順序は全テスト問題サイズで頑健に成立する。第一モーメント法と同一の構造的パラメータ L は、解の存在だけでなく計算コストのスケーリングにも関与する。ただし、その効き方はソルバー依存の感度指数 c によって割り引かれる。

本稿が与えるのは、構造を維持できる状態集合が明示的に定まるドメインにおいて、L という座標がどこまで存在と発見の双方を整理できるかの一つの具体例である。直接の貢献は、存在と発見の分離、感度指数 c の導入、および α-n 非対称性の明示にある。構造持続理論の全体を SAT で証明したと主張するものではない。
\end{document}
